% The edge-edge binning worked through on one scene: (a) the binning, then the
% walk of each box that emits a pair, in the order the cells put them.
%
% No prose inside the drawing. What the panels mean belongs in the caption,
% where it can be read at caption size and does not compete with the marks.
%
% Grids are drawn line by line rather than with TikZ's `grid`, which places its
% lines at multiples of the step measured from the origin and not from the
% corner of the rectangle given to it.
%
% The seven boxes and the grid are real: cells are 0.62 wide, six by four,
% numbered row-major from the bottom left, and every cell index, walk range and
% skipped column below follows from the coordinates in \scenebox.
% Drawn at 0.62 to the cell and scaled up to the text width. Coordinates scale,
% node text and line widths do not, so the panels grow and the labels keep their
% size -- which is why the panel offsets below are coordinates and not lengths.
\begin{tikzpicture}[
  scale=1.22,
  font=\scriptsize,
  lbl/.style={font=\scriptsize, text=sccdInk, inner sep=1.5pt},
  tiny/.style={font=\tiny, text=sccdInk, inner sep=1pt},
  ttl/.style={font=\scriptsize\bfseries, text=sccdInk, inner sep=0pt},
  idle/.style={draw=sccdInk!22, line width=0.5pt},
  query/.style={draw=sccdTeal, line width=1.0pt},
  found/.style={draw=sccdPlum, line width=1.0pt},
  full/.style={fill=sccdTeal!14},
  empty/.style={fill=sccdMoss!30},
  beyond/.style={fill=sccdInk!6},
  % The span one binary search over the row's prefix maximum steps over. The
  % tint is translucent so the cell fills beneath it still read: a searched cell
  % can be covered, empty, or neither.
  search/.style={draw=sccdAmber, line width=0.8pt,
                 dash pattern=on 2pt off 1.4pt,
                 fill=sccdAmber, fill opacity=0.16},
]

% The columns a row's lower-bound search steps over, as c0..c1 on row r.
\newcommand{\searchspan}[3]{%
  \draw[search] ({#1*0.62},{#3*0.62}) rectangle ({(#2+1)*0.62},{(#3+1)*0.62});}

% One cell, addressed by column and row.
\newcommand{\cellfill}[3]{%
  \fill[#3] ({#1*0.62},{#2*0.62}) rectangle ({(#1+1)*0.62},{(#2+1)*0.62});}

% label / x0 / y0 / x1 / y1 -- the minimum corner is (x0,y0).
\def\scenebox{%
  a/0.35/0.30/1.10/1.00, b/1.05/0.20/1.65/0.75, c/0.80/0.85/1.50/1.50,
  d/1.95/0.45/2.60/1.55, e/2.40/1.10/3.10/1.70, f/0.30/1.75/0.90/2.30,
  g/0.75/1.95/1.45/2.40}

% The grid, drawn line by line. #1 is unused; kept as a command for brevity.
\newcommand{\scenegrid}{%
  \foreach \x in {0,0.62,1.24,1.86,2.48,3.10,3.72}
    \draw[draw=sccdInk!30, line width=0.3pt] (\x,0) -- (\x,2.48);
  \foreach \y in {0,0.62,1.24,1.86,2.48}
    \draw[draw=sccdInk!30, line width=0.3pt] (0,\y) -- (3.72,\y);
  \draw[draw=sccdInk!45, line width=0.4pt] (0,0) rectangle (3.72,2.48);}

% Every box in the faint style, then the minimum corner each is binned at.
\newcommand{\sceneidle}{%
  \foreach \n/\xa/\ya/\xb/\yb in \scenebox {
    \draw[idle] (\xa,\ya) rectangle (\xb,\yb);
    \fill[sccdInk!30] (\xa,\ya) circle (1.0pt);}}

% ============================ (a) the binning ================================
\begin{scope}
  \node[ttl, anchor=south west] at (0,2.64) {(a) one cell per box};

  \scenegrid
  \foreach \n/\xa/\ya/\xb/\yb in \scenebox {
    \draw[draw=sccdInk!55, line width=0.7pt] (\xa,\ya) rectangle (\xb,\yb);
    \fill[sccdInk] (\xa,\ya) circle (1.4pt);
    \node[lbl, anchor=south west, inner sep=1pt] at ({\xb+0.02},\yb) {$\n$};}

  % Cell numbers last, so no box can hide one, and in the corner the box labels
  % do not use.
  \foreach \r in {0,1,2,3}
    \foreach \c in {0,1,2,3,4,5} {
      \pgfmathtruncatemacro{\id}{\r*6+\c}
      \node[tiny, text=sccdInk!45, anchor=south west, fill=white,
            inner sep=0.5pt] at ({\c*0.62+0.03},{\r*0.62+0.03}) {\id};}
\end{scope}

% ============================= (b) box a walks ===============================
% a is binned in cell 0 and its maximum corner in cell 7, so it covers cells 0,
% 1, 6 and 7 -- rows 0-1 crossed with columns 0-1. On row 1 the running maximum
% starts the scan at column 1: cell 6 is empty, so nothing there reaches back.
\begin{scope}[shift={(4.32,0)}]
  \node[ttl, anchor=south west] at (0,2.64) {(b) $a$ covers cells $0,1,6,7$};

  \cellfill{0}{0}{full}  \cellfill{1}{0}{full}  \cellfill{1}{1}{full}
  \cellfill{0}{1}{empty}
  \fill[beyond] (1.24,0) rectangle (3.72,1.24);
  \scenegrid
  \sceneidle

  \draw[found] (1.05,0.20) rectangle (1.65,0.75);
  \draw[found] (0.80,0.85) rectangle (1.50,1.50);
  \draw[query] (0.35,0.30) rectangle (1.10,1.00);
  \fill[sccdTeal] (0.35,0.30) circle (1.4pt);
  % Row 1: the search starts the row at column 1, stepping over cell 6.
  \searchspan{0}{0}{1}

  \node[lbl, text=sccdPlum, anchor=south west] at (1.67,0.75) {$b$};
  \node[lbl, text=sccdPlum, anchor=south west] at (1.52,1.50) {$c$};
  \node[lbl, text=sccdTeal, anchor=north east] at (0.33,0.28) {$a$};
\end{scope}

% ============================= (c) box d walks ===============================
% d is binned in cell 3 and its maximum corner in cell 16, so it covers cells 3,
% 4, 9, 10, 15 and 16 -- rows 0-2 crossed with columns 3-4. Row 2 holds only f,
% whose maximum on the first axis is 0.90 and so never reaches d: the running
% maximum returns a column past 4 and the whole row is skipped without being
% read.
\begin{scope}[shift={(0,-3.40)}]
  \node[ttl, anchor=south west] at (0,2.64) {(c) $d$ covers cells $3,4,9,10,15,16$};

  % Only 3 and 9 hold anything: d itself and e. Cells 4 and 10 are read and
  % yield nothing, and 15 and 16 are stepped over by the running maximum.
  \cellfill{3}{0}{full}  \cellfill{3}{1}{full}
  \cellfill{4}{0}{empty} \cellfill{4}{1}{empty}
  \cellfill{3}{2}{empty} \cellfill{4}{2}{empty}
  \fill[beyond] (3.10,0) rectangle (3.72,1.86);
  \scenegrid
  \sceneidle

  \draw[found] (2.40,1.10) rectangle (3.10,1.70);
  \draw[query] (1.95,0.45) rectangle (2.60,1.55);
  \fill[sccdTeal] (1.95,0.45) circle (1.4pt);
  % Row 1: the search steps over columns 0-2, left of $d$'s own columns, where a
  % partner binned lower could still reach back. Row 2: it returns a column past
  % 4 and steps over the whole row.
  \searchspan{0}{2}{1}
  \searchspan{0}{4}{2}

  \node[lbl, text=sccdPlum, anchor=south west] at (3.12,1.70) {$e$};
  \node[lbl, text=sccdTeal, anchor=north east] at (1.93,0.43) {$d$};
\end{scope}

% ============================= (d) box f walks ===============================
% f is binned in cell 12 and its maximum corner in cell 19, so it covers cells
% 12, 13, 18 and 19 -- rows 2-3 crossed with columns 0-1. On row 3 the running
% maximum starts at column 1, cell 18 being empty.
\begin{scope}[shift={(4.32,-3.40)}]
  \node[ttl, anchor=south west] at (0,2.64) {(d) $f$ covers cells $12,13,18,19$};

  \cellfill{0}{2}{full}  \cellfill{1}{3}{full}
  \cellfill{1}{2}{empty} \cellfill{0}{3}{empty}
  \fill[beyond] (1.24,1.24) rectangle (3.72,2.48);
  \scenegrid
  \sceneidle

  \draw[found] (0.75,1.95) rectangle (1.45,2.40);
  \draw[query] (0.30,1.75) rectangle (0.90,2.30);
  \fill[sccdTeal] (0.30,1.75) circle (1.4pt);
  % Row 3: the search starts the row at column 1, stepping over cell 18.
  \searchspan{0}{0}{3}

  \node[lbl, text=sccdPlum, anchor=south west] at (1.47,2.40) {$g$};
  \node[lbl, text=sccdTeal, anchor=north east] at (0.28,1.73) {$f$};
\end{scope}

% ================================== key ======================================
\begin{scope}[shift={(0,-4.30)}]
  \fill[sccdInk] (0.11,0.22) circle (1.4pt);
  \node[lbl, anchor=west] at (0.26,0.22) {the minimum corner a box is binned at};
  \fill[full] (0.00,-0.08) rectangle ++(0.22,0.16);
  \draw[draw=sccdInk!30, line width=0.3pt] (0.00,-0.08) rectangle ++(0.22,0.16);
  \node[lbl, anchor=west] at (0.26,0.00) {a covered cell holding a box};
  \fill[empty] (0.00,-0.38) rectangle ++(0.22,0.16);
  \draw[draw=sccdInk!30, line width=0.3pt] (0.00,-0.38) rectangle ++(0.22,0.16);
  \node[lbl, anchor=west] at (0.26,-0.30) {a covered cell that is empty};
  \fill[beyond] (0.00,-0.68) rectangle ++(0.22,0.16);
  \draw[draw=sccdInk!30, line width=0.3pt] (0.00,-0.68) rectangle ++(0.22,0.16);
  \node[lbl, anchor=west] at (0.26,-0.60) {past the box's own maximum corner};
  \draw[search] (0.00,-0.98) rectangle ++(0.22,0.16);
  \node[lbl, anchor=west] at (0.26,-0.90)
    {the columns one lower-bound search steps over};
\end{scope}
\end{tikzpicture}
